\section{Limitations and Conclusion}
\label{sec:limitations}

\paragraph{Conclusion.}
Conditional unbiasedness is sufficient but not necessary for target-valid CFR
feedback at an explicit finite public cut.  One-sided target debit supplies a
pathwise alternative: it transfers delivered regret to the dense target game,
makes complete coverage a zero-debit recovery action, and supports causal
width and order decisions without future candidate replay.  The equal-metadata
converse identifies the required semantic information, while the two-domain
evidence connects it to a width/order crossover and a held-out dynamic-width
consequence.  Public-feedback granularity can therefore be treated as a
controlled learning decision with an explicit target-game account.

\paragraph{Scope and limitations.}
The theory covers finite additive public cuts, event-faithful alternating
signed RM under uniform weights, and RM+ under a predetermined nondecreasing
output clock.  Dynamic control is evaluated in public Leduc; HUNL establishes
the frontier but not an end-to-end DB-CFR speedup.  The matched-work result
counts installed recurrence components and reports its exact selected-path
audit separately; a full-state asynchronous prototype is limited by snapshot
and memory traffic.  Nested cuts and other learner recurrences require separate
transfer arguments.  Compact within-origin factorization is the next step
toward online DB-CFR in large HUNL endgames.
